\chapter{Volatilidade e Dependência Avançadas}

\section{Além do GARCH}

O Capítulo~31 cobriu GARCH, EGARCH e GJR-GARCH. Este capítulo trata de
volatilidade estocástica, correlações dinâmicas, copulas, wavelets e
processos de pontos auto-excitantes.

\section{Volatilidade Estocástica}

Modelos GARCH modelam a variância condicional como função
determinística de erros quadrados passados. A Volatilidade Estocástica
(Taylor, 1986) trata o log-volatilidade como um processo AR(1) latente:

\[
y_t = \exp(h_t/2) \varepsilon_t, \qquad
h_t = \mu + \phi(h_{t-1} - \mu) + \eta_t
\]

\begin{lstlisting}[caption={Volatilidade Estocástica}]
let n = 500
generate df y = 0.02 + rnormal(n, 0, exp(0.5*rnormal(n, 0, 1)))

let m_sv = sv(df, y, iter=500)
print(m_sv)
\end{lstlisting}

\texttt{sv} reporta $\mu$, $\phi$ (persistência), $\sigma_\eta$
(volatilidade-da-volatilidade) e o caminho de volatilidade condicional.

\section{Correlação Condicional Dinâmica}

DCC-GARCH (Engle, 2002) combina modelos GARCH(1,1) univariados com uma
equação dinâmica para a matriz de correlação condicional. É útil para
modelar dependência variante no tempo entre retornos.

\begin{lstlisting}[caption={DCC-GARCH}]
let m_dcc = dcc_garch(y1 ~ y2, df)
print(m_dcc)
\end{lstlisting}

A saída inclui os parâmetros DCC ($\alpha$, $\beta$), os parâmetros
GARCH univariados, volatilidades condicionais e matrizes de correlação
dinâmica.

\section{Copulas}

\texttt{copula} estima um modelo bivariado de dependência via Teorema de
Sklar. Separa as distribuições marginais da copula que as une.

\begin{lstlisting}[caption={Copula}]
let m_cop = copula(y1 ~ y2, df, type="gaussian")
print(m_cop)
\end{lstlisting}

As copulas disponíveis incluem Gaussiana, Clayton, Gumbel e Frank. A
saída reporta o $\tau$ de Kendall, o $\rho$ de Spearman e o parâmetro
da copula.

\texttt{tv\_copula} estende isso com dependência variante no tempo impulsionada
por uma dinâmica tipo GARCH para o parâmetro da copula.

\begin{lstlisting}[caption={Copula de dependência variante no tempo}]
let m_tvc = tvcopula(y1 ~ y2, df, type="gaussian")
print(m_tvc)
\end{lstlisting}

\texttt{tvcopula} é experimental e exige \texttt{--features experimental}.

\section{Wavelets}

A Transformada Wavelet de Máxima Sobreposição Discreta (MODWT)
decompõe uma série em componentes tempo-frequência sem decimação.

\begin{lstlisting}[caption={MODWT}]
let m_modwt = modwt(df, y, scales=4)
print(m_modwt)
\end{lstlisting}

Útil para isolar componentes cíclicos em diferentes frequências.

\section{Processos de Hawkes}

Um processo de Hawkes é um processo de pontos auto-excitante. Cada
evento aumenta a intensidade condicional de eventos futuros por uma
quantidade que decai exponencialmente.

\begin{lstlisting}[caption={Processo de Hawkes}]
let m_hawkes = hawkes(df, time, T=100)
print(m_hawkes)
\end{lstlisting}

A taxa de ramificação $\eta = \alpha / \beta$ indica quanto cada
evento contribui para eventos futuros. $\eta < 1$ é necessário para
estacionariedade.
